\chapter{Stool changes}
\index{Stool changes}

\medskip
\noindent\textit{Educational reference; always seek professional advice for individual care.}

\section*{Definition and Overview}
Stool changes describe any alteration in the usual appearance, consistency, colour, frequency, or ease of passing a bowel motion. They are a symptom or a sign rather than a diagnosis in their own right, and their meaning depends heavily on context: the sudden onset of loose stools during a tummy bug carries a very different implication from weeks of pencil-thin stools in an older person, or black, tarry motions in someone taking blood-thinning medication. Normal stool is highly variable across individuals. Most people pass anywhere from three stools a day to three a week, and what matters is change from that person's own usual pattern, combined with other symptoms such as pain, bleeding, weight loss, or fever. The stool itself is mostly water and undigested fibre, together with bacteria and waste products, and its appearance is influenced by transit time through the bowel, diet, hydration, bile pigments, and medication. This chapter explains the common causes of stool changes, how to describe them accurately, which features are concerning, how doctors investigate them, and the treatments and self-care measures that usually help. Throughout, the guiding principle is that a single change with a clear cause is usually harmless, while persistent, unexplained, or severe changes warrant medical assessment.

\section*{Epidemiology}
Stool changes are among the most common reasons people consult a general practitioner. Episodes of acute diarrhoea affect the majority of adults and children every few years, most caused by short-lived viral or food-related infections that settle within a few days. Constipation is similarly widespread, affecting perhaps one in six adults at any time, and is especially common in older people, during pregnancy, and in those taking certain medications. Chronic conditions that change the stool, such as irritable bowel syndrome, inflammatory bowel disease, and coeliac disease, together affect a substantial minority of the population, with irritable bowel syndrome alone estimated to affect around one in ten people. Because the lifetime risk of bowel cancer in many countries is among the highest in the world, with roughly one in 20 people developing it during their lifetime, unexplained persistent changes in bowel habit in older adults are taken seriously and investigated promptly. The pattern of stool changes differs with age: toddlers commonly pass loose stools with teething and mild infections, young adults frequently experience stress-related changes, and older adults are more likely to have structural causes such as polyps, cancer, or diverticular disease.

\section*{Aetiology and Pathophysiology}
Stool appearance is set by how long contents spend in the bowel, how much water is absorbed, and what the bowel lining releases or bleeds. When transit is rapid, water absorption is cut short and stools are loose; when transit is slow, more water is absorbed and the stool becomes hard and infrequent. Inflammation, infection, and malabsorption can all disrupt this balance, while bleeding at various points in the gut colours the stool differently.

\begin{itemize}
    \item \textbf{Infection:} viruses such as norovirus and rotavirus, and bacteria such as \textit{Salmonella}, \textit{Campylobacter}, and \textit{E.~coli}, inflame the bowel lining and speed transit, causing watery diarrhoea
    \item \textbf{Diet and lifestyle:} low fibre, inadequate fluid, excessive caffeine or alcohol, and certain foods can cause constipation or looser stools
    \item \textbf{Medications:} antibiotics, laxatives, iron supplements, antacids, opioids, and many others can change stool frequency, colour, or consistency
    \item \textbf{Irritable bowel syndrome:} a common disorder of gut-brain interaction that produces alternating constipation and diarrhoea with bloating and pain
    \item \textbf{Inflammatory bowel disease:} ulcerative colitis and Crohn disease cause chronic inflammation, frequent loose or bloody stools, and urgency
    \item \textbf{Malabsorption:} coeliac disease, pancreatic insufficiency, and lactose intolerance reduce nutrient and water absorption, producing pale, bulky, foul, floating stools
    \item \textbf{Bleeding:} bright red blood usually indicates lower bowel or anal bleeding; black, tarry stools indicate bleeding higher up where blood is digested
    \item \textbf{Structural disease:} polyps, diverticular disease, and colorectal cancer can change calibre and pattern of stools and cause bleeding
    \item \textbf{Endocrine and metabolic:} an overactive thyroid speeds the bowel; diabetes and underactive thyroid can cause constipation
\end{itemize}

\section*{Clinical Features}
Stool changes present in a range of patterns, and describing them accurately is the first step toward a cause.

\begin{itemize}
    \item \textbf{Diarrhoea:} loose or watery stools passed more than three times a day; acute episodes are usually infective, while persistent diarrhoea suggests chronic conditions or malabsorption
    \item \textbf{Constipation:} infrequent, hard, or difficult-to-pass stools, often with straining, bloating, or a feeling of incomplete emptying
    \item \textbf{Blood in the stool:} bright red blood on the paper or stool usually means haemorrhoids or anal fissures, whereas blood mixed through the stool suggests higher bleeding
    \item \textbf{Black, tarry stools:} digested blood from the stomach or upper intestine, often with iron or bismuth products as a harmless alternative
    \item \textbf{Pale or clay-coloured stools:} reduced bile flow, seen with liver or gallbladder problems, sometimes with dark urine and yellowing of the skin
    \item \textbf{Mucus in the stool:} can accompany inflammation, infection, or irritable bowel syndrome, though small amounts are normal
    \item \textbf{Floating or very foul stools:} suggests excess gas or fat, as in malabsorption or pancreatic disease
    \item \textbf{Changed calibre:} persistently narrow or ribbon-like stools can indicate a narrowing in the lower bowel
    \item \textbf{Changed frequency:} an unexplained increase or decrease lasting more than a few weeks
\end{itemize}

\section*{Differential Diagnosis}
Because stool changes have many causes, the clinician works through the most likely and most serious possibilities first.

\begin{itemize}
    \item \textbf{Infectious gastroenteritis:} sudden loose stools with nausea, vomiting, or fever; settle within days and often follow exposure to food, travel, or contacts
    \item \textbf{Irritable bowel syndrome:} chronic, crampy pain with altered stool pattern and bloating; pain eases with defecation and there is no blood or weight loss
    \item \textbf{Inflammatory bowel disease:} persistent diarrhoea with blood, mucus, urgency, weight loss, and anaemia; diagnosed by colonoscopy and biopsies
    \item \textbf{Coeliac disease:} chronic loose stools, bloating, fatigue, and poor growth in children; confirmed by blood tests and bowel biopsy
    \item \textbf{Colorectal cancer:} new, persistent change in bowel habit, especially after age 50, with blood, anaemia, or weight loss; must be excluded with colonoscopy
    \item \textbf{Drug-related changes:} a careful medication history often reveals antibiotics, iron, opioids, or laxative misuse as the culprit
    \item \textbf{Haemorrhoids and fissures:} bright red bleeding and discomfort with otherwise normal stools
    \item \textbf{Diverticular disease:} lower abdominal pain with changes in bowel habit, more common over age 60
\end{itemize}

\section*{When to Seek Medical Care}
A single loose or hard stool needs no treatment, but medical review is sensible when changes persist beyond a few days or are accompanied by worrying features. See a general practitioner if stool changes last more than three weeks, if they are associated with unexplained weight loss, persistent abdominal pain, anaemia symptoms such as fatigue and breathlessness, or blood mixed through the stool, or if they occur alongside fever, night sweats, or poor appetite. Older adults, people with a family history of bowel cancer, and those with inflammatory bowel disease or coeliac disease should have persistent changes assessed early. Infants and young children who develop diarrhoea or vomiting should be assessed quickly if they show signs of dehydration, such as a dry mouth, fewer wet nappies, listlessness, or sunken eyes.

\textbf{Contact your local emergency services for:}
\begin{itemize}
    \item Large amounts of blood, especially with light-headedness, collapse, or a racing pulse
    \item Severe abdominal pain that is worsening, with a swollen or rigid abdomen
    \item Signs of severe dehydration or shock, such as confusion, fainting, or very little urine
    \item Bloody diarrhoea with high fever, particularly in a young child or older person
\end{itemize}

\section*{Diagnosis and Investigations}
Most acute stool changes need no tests and settle with supportive care. When they persist or raise concern, investigation is targeted.

\begin{itemize}
    \item \textbf{History:} the clinician asks about onset, pattern, stool appearance, diet, travel, medicines, family history, and associated symptoms
    \item \textbf{Physical examination:} the abdomen is felt for tenderness or masses, and a rectal examination may be performed
    \item \textbf{Stool tests:} samples can be checked for infection, blood, or markers of inflammation such as calprotectin
    \item \textbf{Blood tests:} full blood count for anaemia, and checks for coeliac antibodies, thyroid, and liver function
    \item \textbf{Colonoscopy:} a camera examination of the large bowel, the definitive test for inflammation, polyps, and cancer
    \item \textbf{Flexible sigmoidoscopy:} examination of the lower bowel only, sometimes used for rectal bleeding
    \item \textbf{Imaging:} CT scans are used when a structural problem or complication is suspected
    \item \textbf{Stool-blood screening:} screening programmes may offer home testing at defined ages and intervals to detect hidden blood; local eligibility varies
\end{itemize}

\section*{Management}
Treatment is directed at the underlying cause, and many stool changes respond to simple measures.

\begin{itemize}
    \item \textbf{Fluid replacement:} drink plenty of water; for significant diarrhoea, oral rehydration solutions replace salt and sugar lost from the body
    \item \textbf{Diet for diarrhoea:} plain foods such as rice, toast, and bananas are well tolerated, while spicy, rich, or fatty foods may worsen symptoms
    \item \textbf{Diet for constipation:} increase fibre gradually from fruit, vegetables, and whole grains, and drink more fluid
    \item \textbf{Laxatives:} bulking agents, stool softeners, or stimulant laxatives, used short term under advice for troublesome constipation
    \item \textbf{Antidiarrhoeal medicines:} agents such as loperamide can reduce loose stools but should be avoided in suspected serious infection
    \item \textbf{Antibiotics:} only for confirmed bacterial infections or specific conditions; most gut infections do not need them
    \item \textbf{Medication review:} adjusting or ceasing a culprit drug may resolve the change
    \item \textbf{Treating the underlying disease:} corticosteroids and immune-suppressing drugs for inflammatory bowel disease, a gluten-free diet for coeliac disease, and surgical treatment for cancer or polyps
\end{itemize}

\section*{Complications}
Most stool changes resolve without harm, but some patterns carry risk if not addressed.

\begin{itemize}
    \item \textbf{Dehydration:} the main danger of acute diarrhoea, especially in infants, the elderly, and those with chronic illness
    \item \textbf{Anaemia:} chronic hidden bleeding causes iron deficiency, fatigue, and reduced exercise capacity
    \item \textbf{Bowel obstruction:} unrelieved constipation can progress to faecal impaction and blockage, particularly in older people
    \item \textbf{Perforation and peritonitis:} severe inflammatory or structural bowel disease can rarely cause a tear and infection of the abdominal cavity
    \item \textbf{Delayed cancer diagnosis:} the most serious consequence of ignoring persistent stool changes, which is why timely investigation matters
    \item \textbf{Nutritional deficiency:} long-term malabsorption leads to weight loss and deficiency of vitamins and minerals
\end{itemize}

\section*{Prognosis and Outcomes}
The outlook for stool changes depends entirely on the cause. Acute infective diarrhoea settles in most people within a week without treatment, and constipation responds well to dietary change and simple measures in the great majority of cases. Chronic conditions such as irritable bowel syndrome are managed long term, with most people improving substantially with dietary advice, stress management, and appropriate medications, even though symptoms may come and go. Inflammatory bowel disease and coeliac disease are lifelong conditions, but modern treatment and strict diet allow most people to lead full lives with good symptom control. When investigation reveals polyps or early bowel cancer, removal or surgery is often curative, which is why persistent changes should never be ignored. In many countries, screening programs have improved early detection markedly. The most important predictors of a good outcome are early assessment, an accurate diagnosis, and treatment of the underlying condition rather than repeated self-treatment of symptoms.

\section*{Prevention and Self-Care}
\begin{itemize}
    \item Keep a food and symptom diary to identify triggers such as lactose, spicy foods, or caffeine
    \item Eat a high-fibre diet with plenty of fruit, vegetables, legumes, and whole grains, increasing fibre gradually
    \item Drink adequate fluid throughout the day to keep stools soft
    \item Wash hands well after the toilet and before preparing food to prevent infectious diarrhoea
    \item Practise food hygiene with proper refrigeration, cooking, and separation of raw meats
    \item Take part in the bowel-cancer screening programme available where you live when invited
    \item Review medicines with a pharmacist or doctor if they are causing digestive symptoms
    \item Seek medical assessment rather than repeatedly self-treating unexplained changes
\end{itemize}

\section*{Sources and Further Reading}
The following organisations provide reliable information about stool changes, bowel symptoms, and digestive health.

\begin{itemize}
    \item NHS. \textit{Diarrhoea and constipation: symptoms and treatment.} https://www.nhs.uk/
    \item Mayo Clinic. \textit{Blood in stool and changes in bowel habits.} https://www.mayoclinic.org/
    \item healthdirect Australia. \textit{Bowel habit changes and faecal occult blood testing.} https://www.healthdirect.gov.au/
    \item Australian Department of Health and Aged Care. \textit{National Bowel Cancer Screening Program.} https://www.health.gov.au/
    \item World Health Organization (WHO). \textit{Diarrhoeal disease: fact sheet.} https://www.who.int/
    \item GESA (Gastroenterological Society of Australia). \textit{Information for patients about bowel symptoms.} https://www.gesa.org.au/
\end{itemize}

\clearpage
